\documentclass[a4paper,12pt]{memoir}
\usepackage{graphicx}
\usepackage[usenames,svgnames]{xcolor}
\chapterstyle{bianchi}
\usepackage{a4wide,amsfonts,appendix}
\usepackage[show]{ed}
\usepackage{url}

% Less detailed TOC
\setcounter{tocdepth}{1}


\begin{document}
%% Some shades
\definecolor{Dark}{gray}{0.2}
\definecolor{MedDark}{gray}{0.4}
\definecolor{Medium}{gray}{0.6}
\definecolor{Light}{gray}{0.8}
\newlength{\drop}% for my convenience
\newlength{\dropt}% for my convenience

\renewcommand{\maketitle}{\begingroup%
\drop = 0.1\textheight
\dropt = 0.05\textheight
\fboxsep 0.5\baselineskip
\sffamily
\vspace*{\dropt}
\centering
{\textcolor{NavyBlue}{\HUGE The Structure of Mathematical\\[3mm] Expressions}}\par
\vspace{0.5\drop}
\fcolorbox{NavyBlue}{Dark}{\textcolor{white}{\normalfont\Large
An {\textsc{arXiv}} Case Study}}\par
\vspace{\drop}
{\Large Deyan Ginev and Bruce R. Miller}\par
\vspace{0.05\textheight}
{\large National Institute of Standards and Techonology}\par
\vfill
{\begin{center}\today\end{center}}
\vspace*{\drop}
{\hfill\includegraphics{NIST}}
\endgroup}
\makeatother

\pagestyle{empty}
\maketitle
\newpage
\pagestyle{plain}
% Remove parskip for toc
\setlength{\parskip}{0ex plus 0.5ex minus 0.2ex}
\tableofcontents

\chapter{Introduction}
In this study, we survey the notational diversity of present-day mathematical expressions, in order to uncover its linguistic phenomena.
\section{Motivation}
We want to enable machine-reading of formulas, in order to provide a variety of user-assistance services, such as semantic search, text-to-speech synthesis, semantic interactions (definition lookup), as well as computer algebra support (``evaluate subexpressions on demand'') and ultimately computer verification (``does that proof step really hold?'').\ednote{expand}
\section{Related Resources}
Notation census, beginnings of study are in Deyan's thesis, Naproche and FMathL have examples, but no real systematic study.\ednote{expand}

\section{Experimental Setup}
The primary corpus on which we base this investigation is the Cornel pre-print archive ``arXiv''\ednote{cite here}, consisting of over 700,000 articles in 37 scientific subfields.

As a secondary resource, we we will also consult entry-level literature on highschool mathematics, in order to exhibit basic phenomena, as well as to demonstrate phenomena apriori known to the authors

\chapter{A Study of Mathematical Syntax}

\section{Basics}
\subsection{Foundations}

\ednote{arithmetic, grouping fences and equality}
\ednote{basic relations and orderings}
\ednote{arithmetic and algebraic sequences?} 

% Maybe add a high school subsection/subsubsection:
\subsection{High School}
\ednote{geometry here, otherwise a separate geometry subsection}
\ednote{trigonometry, complex and rational numbers}

\section{Discrete math}
\subsection{Set Theoretic Notations}
\ednote{elementhood, inclusions, set constructors, overloaded arith ops}
\ednote{also maps : domains -> codomains, xRy notations}

\begin{entries}{Set Theory Notations, Part 1}{study:sets1}
\entry{$ W \in \CP \cap \CZ $}{set membership}%
{ \MathTree [.\is{infix,relation}{$\in$} W [.$\cap$ $\CP$ $\CZ$ ] ] }%
{set ops precede set relations, [Train1]}
% \entry{$ \; \nu : \overset{n}{\times} \BV \ra \BR$}{a map}{typed modifier}{$n$-ary cross-product, [Train1]}
% \entry{\footnotesize $\CZ^* = \{ X \in \CV \mid \om(X,W) \in \BZ, \text{ for all }  W \in \CZ \}$}{definition to set}{typed modifier}{NL mixins, quantified relation, [Train1]}
% \entry{$\text{span}_{\BR} \{ W_{1} , \dots ,W_{g} \} $}{span of a set}{prefix operator modified}{set operators can take fenced yet not simply \textit{grouped} arguments, [Train1]}

% \entry{$\mathcal{BS}_{\CP} = \underset{\bq \in (\BZ/ k \BZ)^g}{\cup} \, \La_{\bq}$}{definitional assignment}{infix relation}{$n$-ary union, ranges over subscript, [Train1]}
% \entry{ $L_{12} = L_{\CP_1} \cap L_{\CP_2} \neq 0$}{non-transversal intersection}{modified chained}{acts as relation, is it modified or chained?, [Train1]}
% \entry{ $\al \in GL(g,\BZ) \subset Sp(2g,\BZ)$}{statement}{infix relation}{membership comes after set inclusion, due to well-typedness!, [Train1]}
\end{entries}

%\Tree [.\is{infix,relation}{$\in$} [W .$\cap$ [$\CP$ $\CZ$] ] ]
\subsection{Logical Operators}
\ednote{classic logic, HOL, type theories} 
\subsection{Combinatorics}
\ednote{Infinite sums}
\ednote{binomials, combinations, permutations,}
\subsection{Number Theory}
\ednote{modulo modifiers}
\ednote{tuples}
\ednote{divisibility notations $a \mid b $ and $b/a$}
\ednote{DLMF sneaky notations}
\subsection{Graph Theory}
\ednote{edge and vertex notations}
\ednote{incidence and adjacency notations}
\ednote{Wiki is very nice: \url{http://en.wikipedia.org/wiki/Glossary\_of\_graph\_theory}}
\subsection{Algebra}
\ednote{vectors}
\ednote{maps and complements}
\ednote{groups}
\ednote{lattices}
\documentclass{article}
\usepackage{latexml}
\lxDeclare[role=ID]{$a$}%
\lxDeclare[role=ID]{$b$}%
\lxDeclare[role=ID]{$c$}%
\lxDeclare[role=ID]{$d$}%
\lxDeclare[role=ID]{$e$}%
\lxDeclare[role=FUNCTION]{$f$}%
\lxDeclare[role=FUNCTION]{$g$}%
\lxDeclare[role=FUNCTION]{$h$}%
\lxDeclare[role=OPFUNCTION]{$F$}%
\lxDeclare[role=OPFUNCTION]{$G$}%
\lxDeclare[role=OPFUNCTION]{$H$}%
\begin{document}
\section{Basic Functions}
\[ g(a) = g(b) \]
\[ f(a) = g(h(c)) \]
\[ h(c) = \sin a^2 \]
\[ \sin f(a) \]

\section{Trigonometric functions}
\[ \sin a + \tan b \]
\[ \sin a \cos b \neq \sin(a\cos b)\]
% Parses wrong ? but how to clarify?
\[ \sin 2\pi \neq \log 2x^2 \]

\section{Missing arguments}
% Presumably these are products
\[ f g h \]
\[ f g h a \]
\section{Operator-like functions}
% But these are applied!
\[ F G H \]
\[ F G H a \]

\[ f'(a) = a^2 \]
\[ f^2(a) = g^b(2) \]

\section{Scripts}
\[ \ln a = \log_e a = \ln_{+} b\]
\end{document}


\section{Continuous math}
\subsection{Calculus}
\ednote{differentials, integrals, limits, remember brownian motion integral notations!}
\subsection{Probability}

\ednote{Bayes formula with multiple denotations of P} % syntactically irrelevant?
\ednote{Various conditional and joint probability notations}

\subsection{Interval Notation and Arithmetic}
\ednote{introduce interval notations, then move to interval arithmetic}
\subsection{Topology}
\ednote{manifold constructors and notations}

\section{Other fields}
\subsection{Quantum Physics}
\ednote{Bra-ket notation}
\ednote{computer science, biology, chemistry...}
$\vdots$

\chapter{Discussion}

\chapter{Conclusion}

\end{document}